\exercise{4.1}{Three upper bounds and a common analysis of the sets}

\goal For each of the twenty-three sets in parts (a)--(w), give
three distinct real upper bounds when such bounds exist. Otherwise,
give the requested answer ``NBA,'' meaning ``not bounded above.''

\idea Exercises 4.2, 4.3, and 4.4 ask about the very same sets.
We first analyze each set once, including the sharpness of its endpoint
bounds. Four answer lists then state the answer to every part of every
exercise and link it to its proof. This keeps the proofs of the
intersections, the prime-number assertion, and the trigonometric values
available without repeating them four times.

\subsubsection*{Definitions and tools for all four exercises}

An \emph{upper bound} $U\in\R$ for a set $S$ must satisfy $s\leq U$
for every $s\in S$. A \emph{lower bound} $L\in\R$ must satisfy
$L\leq s$ for every $s\in S$. Thus, once $U$ is known to be an upper
bound, every number at least $U$ is another upper bound: if $s\leq U$
and $U\leq V$, then $s\leq V$. Similarly, every number at most a known
lower bound is another lower bound.

For a proposed supremum $u$, we prove two things:
\begin{enumerate}
 \item $s\leq u$ for every member $s$ of the set;
 \item for each real $U<u$, there is a member $s>U$.
\end{enumerate}
The second condition means that no number smaller than $u$ is an
upper bound, so $u$ is the \emph{least} upper bound. For a proposed
infimum $\ell$, the corresponding conditions are that $\ell\leq s$
for every member, and that every real $L>\ell$ has some member $s<L$.
If a bound belongs to the set, that member itself proves sharpness:
any smaller upper bound, or larger lower bound, fails on that member.
Thus a maximum is a supremum, and a minimum is an infimum.

All suprema and infima in these four exercises are \emph{real}
numbers. An unbounded-above set has ``NO sup,'' and an unbounded-below
set has ``NO inf,'' in the conventions of this section. These answers
are not assignments of extended-real values.

\step{Tool 1: The Archimedean property and small reciprocals.}
The Archimedean property implies that, for every real $B$, there is
$n\in\N$ with $n>B$. If $B\leq0$, $n=1$ works; if $B>0$, apply
the Archimedean property with $a=1$ and $b=B$.
Consequently, for every $\delta>0$, there is an $n\in\N$ such that
\[
 n>\frac1\delta
 \quad\hbox{and therefore}\quad
 0<\frac1n<\delta.
\]
To justify the last strict inequality, multiply $n>1/\delta$ by
the positive number $\delta$, obtaining $n\delta>1$, and then
divide by the positive integer $n$.

\step{Tool 2: Rational points in an open interval make its endpoints sharp.}
Let $c,d\in\R$ with $c<d$, and consider
\[
 T=\{r\in\Q:c<r<d\}.
\]
The density of $\Q$ shows that $T$ is nonempty. Its defining
inequalities give the upper bound $d$ and the lower bound $c$.
If $U<d$, put $v=\max\{U,c\}$. Both $U$ and $c$ are less than
$d$, so $v<d$. Density provides a rational $r$ with
\[
 \max\{U,c\}<r<d.
\]
Then $c<r<d$, so $r\in T$, and $r>U$, so $U$ is not an upper
bound. This proves $\sup T=d$.

If $L>c$, put $w=\min\{L,d\}$. Both $L$ and $d$ are greater than
$c$, so $c<w$. Density provides a rational $r$ with
\[
 c<r<\min\{L,d\}.
\]
This $r$ belongs to $T$ and satisfies $r<L$, proving that $L$ is
not a lower bound. Hence $\inf T=c$. Neither endpoint belongs to
$T$, since membership requires strict inequalities. The same witnesses
also prove that the real interval $(c,d)$ has supremum $d$ and
infimum $c$, because those rational witnesses are real members of it.

\step{Tool 3: Converting a quadratic inequality into an interval.}
For $c>0$ and $x\in\R$,
\[
 x^2<c^2\quad\Longleftrightarrow\quad -c<x<c.
\]
For the reverse direction, if $-c<x<c$, then $c+x>0$ and
$c-x>0$. Their product is positive, so
\[
 0<(c-x)(c+x)=c^2-x^2,
\]
which says $x^2<c^2$. For the forward direction, if $x\geq c>0$,
then $x^2\geq c^2$, contradicting the strict inequality. If
$x\leq-c$, then $-x\geq c>0$ and
$x^2=(-x)^2\geq c^2$, again a contradiction. Thus neither
$x\geq c$ nor $x\leq-c$ is possible, proving $-c<x<c$.

\step{Tool 4: Named numerical bounds used for the familiar constants.}
We use the familiar estimates
\[
 3<\pi<\frac{22}{7},\qquad 2<e<3.
\]
In particular $2<e<3<\pi<4$, because $22/7<4$ follows from
$22<28$. These estimates suffice to order the finite sets involving
$\pi$ and $e$; decimal approximations are unnecessary.

\subsubsection*{A proof for every set (a)--(w)}

\phantomsection\label{bounds:a}
\step{Analysis (a): The closed interval $[0,1]$.}
Every member $s$ satisfies $0\leq s\leq1$, by the definition of this
interval. Both $0$ and $1$ belong to it. Thus $0$ is its minimum and
infimum, and $1$ is its maximum and supremum. In particular $1,2,3$
are upper bounds and $0,-1,-2$ are lower bounds.

\phantomsection\label{bounds:b}
\step{Analysis (b): The open interval $(0,1)$.}
Every member satisfies $0<s<1$, so $0$ is a lower bound and $1$ is
an upper bound. Tool 2 with $c=0$, $d=1$ proves that no upper bound
smaller than $1$, and no lower bound larger than $0$, can work.
Therefore $\sup(0,1)=1$ and $\inf(0,1)=0$. Neither endpoint is a
member, so neither extremal bound is attained. Upper bounds include
$1,2,3$; lower bounds include $0,-1,-2$.

\phantomsection\label{bounds:c}
\step{Analysis (c): The finite set $\{2,7\}$.}
Its two members satisfy $2\leq2\leq7$ and $2\leq7\leq7$.
Consequently $2$ bounds every member from below and $7$ bounds every
member from above. Each is itself in the set. Hence the minimum and
infimum are $2$, while the maximum and supremum are $7$.
Examples of upper bounds are $7,8,9$; examples of lower bounds are
$2,1,0$.

\phantomsection\label{bounds:d}
\step{Analysis (d): The finite set $\{\pi,e\}$.}
Tool 4 gives $2<e<3<\pi<4$, and therefore $e<\pi$.
Thus every member of this two-element set is at least $e$ and at
most $\pi$. Both endpoints are members, so its infimum is $e$ and
its supremum is $\pi$, both attained. Since every member is less
than $4$, the three distinct numbers $4,5,6$ are upper bounds.
Since every member is greater than $2$, the distinct numbers $2,1,0$
are lower bounds.

\phantomsection\label{bounds:e}
\step{Analysis (e): The reciprocal set $\{1/n:n\in\N\}$.}
For $n\geq1$, positivity and $n\geq1$ give
\[
 0<\frac1n\leq1.
\]
At $n=1$, the value is $1$, so $1$ is the maximum and supremum.
The number $0$ is a lower bound, but it is not a member since a
positive reciprocal cannot equal $0$. If $L>0$, Tool 1 supplies
$n$ with $1/n<L$. This member defeats $L$ as a lower bound, so
$0$ is the greatest lower bound. Thus the infimum is $0$ and is
not attained. Upper bounds include $1,2,3$; lower bounds include
$0,-1,-2$.

\phantomsection\label{bounds:f}
\step{Analysis (f): The singleton $\{0\}$.}
Its only member is $0$. Therefore $0$ is both a lower and an upper
bound, and is both the minimum and the maximum. It follows that
the supremum and infimum both equal $0$. The numbers $0,1,2$ are
three distinct upper bounds, and $0,-1,-2$ are three distinct lower
bounds.

\phantomsection\label{bounds:g}
\step{Analysis (g): The union $[0,1]\cup[2,3]$.}
A member $s$ belongs to at least one of the two intervals.
If $s\in[0,1]$, then $0\leq s\leq1\leq3$. If $s\in[2,3]$,
then $0\leq2\leq s\leq3$. Thus every member satisfies
$0\leq s\leq3$. The member $0$ comes from the first interval,
and the member $3$ comes from the second. Hence the infimum is
the attained minimum $0$, and the supremum is the attained maximum
$3$. The gap between the two intervals changes neither endpoint
bound. Upper bounds include $3,4,5$ and lower bounds include
$0,-1,-2$.

\phantomsection\label{bounds:h}
\step{Analysis (h): The infinite union $\bigcup_{n=1}^{\infty}[2n,2n+1]$.}
Membership in the union means membership in \emph{at least one}
interval. If $s$ belongs to it, there is an $n\in\N$ with
$2n\leq s\leq2n+1$. Because $n\geq1$, we obtain
$2\leq2n\leq s$. The number $2$ belongs to the interval at
$n=1$, namely $[2,3]$. Thus $2$ is the minimum and infimum,
and $2,1,0$ are lower bounds.

There is no upper bound. To see this for an arbitrary proposed
$M\in\R$, use Tool 1 to choose
\[
 n>\max\{1,M/2\}.
\]
The member $2n\in[2n,2n+1]$ satisfies $2n>M$. Thus every proposed
real upper bound fails. This set is NBA and has NO sup.

\phantomsection\label{bounds:i}
\step{Analysis (i): The intersection $\bigcap_{n=1}^{\infty}[-1/n,1+1/n]$.}
Call this set $I$. We prove the exact identity $I=[0,1]$ by both
inclusions. If $0\leq x\leq1$, then for every $n\in\N$,
\[
 -\frac1n<0\leq x\leq1<1+\frac1n.
\]
Thus $x$ belongs to every displayed closed interval, proving
$[0,1]\subseteq I$.

For the other inclusion, suppose $x\notin[0,1]$. If $x<0$, let
$\delta=-x>0$. Tool 1 gives an $n$ with $1/n<\delta=-x$.
Multiplying this inequality by $-1$ reverses its direction:
\[
 -\frac1n>x.
\]
Therefore $x$ fails the lower-endpoint requirement in that particular
interval, so $x\notin I$. If $x>1$, let $\delta=x-1>0$ and
choose $n$ with $1/n<x-1$. Then $1+1/n<x$, so $x$ fails the
upper-endpoint requirement and again is not in $I$. These two cases
exclude every real number outside $[0,1]$, proving $I\subseteq[0,1]$.

It follows that $I=[0,1]$. Its infimum is the member $0$ and its
supremum is the member $1$. Upper bounds are $1,2,3$, for example,
and lower bounds are $0,-1,-2$.

\phantomsection\label{bounds:j}
\step{Analysis (j): The set $\{1-1/3^n:n\in\N\}$.}
The denominator is $3^n$, so the expression is $1-3^{-n}$.
For $n\geq1$, $3^n\geq3$, and hence
\[
 0<\frac1{3^n}\leq\frac13,
 \qquad
 \frac23=1-\frac13\leq1-\frac1{3^n}<1.
\]
At $n=1$, the lower value $2/3$ occurs. Thus $2/3$ is the minimum
and infimum. The number $1$ is an upper bound but is not a member.

To prove that this upper bound is sharp without assuming convergence,
we first verify $3^n\geq n$ for $n\in\N$. At $n=1$, $3\geq1$.
If $3^n\geq n$, multiplication by $3>0$ gives
$3^{n+1}\geq3n$. Also
$3n-(n+1)=2n-1\geq1>0$, so $3n\geq n+1$.
Induction proves the auxiliary inequality.

Now fix a proposed smaller bound $U<1$ and set $\delta=1-U>0$.
Choose $n$ by Tool 1 so that $1/n<\delta$. Since $3^n\geq n>0$,
\[
 \frac1{3^n}\leq\frac1n<1-U,
 \qquad
 1-\frac1{3^n}>U.
\]
This member exceeds $U$, so $U$ is not an upper bound. Therefore
the supremum is $1$, not attained. Upper bounds include $1,2,3$;
lower bounds include $2/3,0,-1$.

\phantomsection\label{bounds:k}
\step{Analysis (k): The set $\{n+(-1)^n/n:n\in\N\}$.}
The numerator $(-1)^n$ equals either $1$ or $-1$. Therefore,
for each $n\in\N$,
\[
 n+\frac{(-1)^n}{n}
 \geq n-\frac1n
 \geq n-1
 \geq0.
\]
Here the first inequality uses $(-1)^n\geq-1$ and division by
$n>0$; the second uses $1/n\leq1$; the third uses $n\geq1$.
At $n=1$, the expression is $1+(-1)/1=0$. Thus $0$ is the
minimum and infimum, and $0,-1,-2$ are lower bounds.

For any proposed real upper bound $M$, choose $n\in\N$ with
$n>M+1$, using Tool 1. The same estimate gives
\[
 n+\frac{(-1)^n}{n}\geq n-1>M.
\]
Thus the set is NBA and has NO sup. The oscillating correction
$(-1)^n/n$ does not keep the growing term $n$ bounded.

\phantomsection\label{bounds:l}
\step{Analysis (l): The rational half-line $\{r\in\Q:r<2\}$.}
All members are less than $2$, so $2$ is an upper bound. If $U<2$,
density provides a rational $r$ with $U<r<2$. That $r$ is a member
and is larger than $U$, so no $U<2$ is an upper bound. Hence the
supremum is $2$. It is not attained because $r=2$ fails the strict
membership inequality. Upper bounds include $2,3,4$.

There is no lower bound: for a proposed $L\in\R$, choose
$n\in\N$ with $n>\max\{1,-L\}$. Then $-n$ is an integer,
hence rational, and $-n<0<2$, so $-n$ is in the set. Also
$n>-L$ gives $-n<L$. Thus every proposed lower bound fails.
The answers for lower boundedness and infimum are NBB and NO inf.

\phantomsection\label{bounds:m}
\step{Analysis (m): The rational set $\{r\in\Q:r^2<4\}$.}
Tool 3 with $c=2$ gives
\[
 \{r\in\Q:r^2<4\}=\{r\in\Q:-2<r<2\}.
\]
Thus $-2$ is a lower bound and $2$ is an upper bound. To see
sharpness explicitly, for $U<2$ density supplies a rational
$\max\{U,-2\}<r<2$, a member exceeding $U$. For $L>-2$ it
supplies a rational $-2<r<\min\{L,2\}$, a member below $L$.
These are exactly the two endpoint arguments of Tool 2. Thus
the supremum is $2$ and the infimum is $-2$.
Neither is a member: $(2)^2=(-2)^2=4$, whereas membership requires
a square strictly less than $4$. Three upper bounds are $2,3,4$;
three lower bounds are $-2,-3,-4$.

\phantomsection\label{bounds:n}
\step{Analysis (n): The rational set $\{r\in\Q:r^2<2\}$.}
Let $c=\sqrt2$, the positive real number with $c^2=2$.
Tool 3 gives
\[
 \{r\in\Q:r^2<2\}=\{r\in\Q:-\sqrt2<r<\sqrt2\}.
\]
For $U<\sqrt2$, density gives a rational
$\max\{U,-\sqrt2\}<r<\sqrt2$; it is a member above $U$.
For $L>-\sqrt2$, density gives a rational
$-\sqrt2<r<\min\{L,\sqrt2\}$; it is a member below $L$.
Together with the bounding inequalities, this proves
\[
 \sup S=\sqrt2,\qquad \inf S=-\sqrt2.
\]
Neither endpoint belongs to the set, since its square equals $2$
instead of being less than $2$. This exclusion already follows from
the strict inequality; it does not require a separate irrationality
argument.

Since $0<\sqrt2<2$, every member lies strictly between $-2$ and
$2$. For the last comparison, $\sqrt2\geq2$ would imply
$2=(\sqrt2)^2\geq4$, a contradiction. Therefore $2,3,4$ are
upper bounds and $-2,-3,-4$ are lower bounds.

\phantomsection\label{bounds:o}
\step{Analysis (o): The negative half-line $\{x\in\R:x<0\}$.}
The number $0$ is an upper bound. If $U<0$, then
\[
 U<\frac U2<0,
\]
because $U/2-U=-U/2>0$ and $U/2$ is still negative.
Thus $U/2$ is a member above the proposed bound $U$, proving
that the supremum is $0$. The value $0$ is not a member.
Upper bounds include $0,1,2$.

For an arbitrary proposed lower bound $L\in\R$, take
$x=\min\{L,0\}-1$. Then
\[
 x<\min\{L,0\}\leq L,
 \qquad x<\min\{L,0\}\leq0.
\]
Hence $x$ is a negative real number below $L$. This proves NBB
and therefore NO inf.

\phantomsection\label{bounds:p}
\step{Analysis (p): The finite set $\{1,\pi/3,\pi^2,10\}$.}
The third entry is $\pi^2$, not a cube. We verify the ordering
using Tool 4. Dividing $3<\pi<22/7$ by the positive number $3$
gives
\[
 1<\frac\pi3<\frac{22}{21}<2,
\]
where $22/21<2$ follows from $22<42$. Squaring the positive
inequalities $3<\pi<22/7$ gives
\[
 9<\pi^2<\frac{484}{49}<10.
\]
The final strict inequality holds since $484<490=10\cdot49$.
Combining these comparisons with $2<9$, we obtain
\[
 1<\frac\pi3<2<9<\pi^2<10.
\]
Thus $1$ is the member of smallest value, and $10$ is the member
of largest value. The infimum is the attained minimum $1$, and
the supremum is the attained maximum $10$. Upper bounds include
$10,11,12$; lower bounds include $1,0,-1$.

\phantomsection\label{bounds:q}
\step{Analysis (q): The finite set $\{0,1,2,4,8,16\}$.}
The displayed members satisfy
\[
 0<1<2<4<8<16.
\]
Every member is at least $0$ and at most $16$, and both $0$ and
$16$ appear in the set. Hence the minimum and infimum equal $0$,
and the maximum and supremum equal $16$. Upper bounds include
$16,17,18$, while lower bounds include $0,-1,-2$.

\phantomsection\label{bounds:r}
\step{Analysis (r): The intersection $\bigcap_{n=1}^{\infty}(1-1/n,1+1/n)$.}
Let $J$ denote this intersection. For each $n\in\N$, $1/n>0$,
so $1-1/n<1<1+1/n$. Therefore $1$ belongs to every one of
the open intervals and hence to $J$.

Now suppose $x\ne1$. If $x<1$, take $\delta=1-x>0$ and use
Tool 1 to choose $n$ with $1/n<1-x$. Rearranging gives
$x<1-1/n$. Thus $x$ lies below the lower endpoint of that
particular open interval and cannot belong to $J$.
If $x>1$, choose $n$ with $1/n<x-1$. Then $1+1/n<x$, and
$x$ lies above the upper endpoint of that interval, so again
$x\notin J$. These two cases exclude every $x\ne1$.

It follows that $J=\{1\}$. Thus both its infimum and its supremum
are $1$, and both are attained. The fact that the intersected
intervals are open does not prevent their intersection from being a
singleton. Three upper bounds are $1,2,3$; three lower bounds are
$1,0,-1$.

\phantomsection\label{bounds:s}
\step{Analysis (s): Reciprocals of positive primes.}
The set is
\[
 P=\{1/p:p\in\N\text{ and }p\text{ is prime}\}.
\]
Every positive prime $p$ satisfies $p\geq2$, so
$0<1/p\leq1/2$. The integer $2$ is prime: its only positive
integer factors are $1$ and $2$. Consequently $1/2$ belongs to
$P$, and is its maximum and supremum. The number $0$ is a lower
bound and is not a member. To show it is the \emph{greatest}
lower bound, we need primes as large as we wish.

\textbf{The needed prime-number fact, with proof.}
First, each integer $m\geq2$ has a prime divisor. The finite
nonempty set
\[
 D=\{d\in\{2,3,\ldots,m\}:d\text{ divides }m\}
\]
contains $m$, so it has a smallest member $p$. If $p$ were not
prime, it would have a positive divisor $a$ other than $1$ and
$p$, with $2\leq a<p$. Writing $p=ab$ and $m=pc$ with integers
$b,c$ gives $m=a(bc)$, so $a$ would divide $m$ and belong to
$D$. That contradicts the minimality of $p$. Hence $p$ is prime.

Now let $B\in\R$ be arbitrary. By Tool 1 choose an integer
$N>\max\{B,2\}$. The integer $m=N!+1$ is at least $2$, so it
has a prime divisor $p$. If $p\leq N$, then $p$ is one of the
factors in the product $N!=1\cdot2\cdots N$. Thus both $N!$ and
$N!+1$ would be divisible by $p$. Subtracting their integer-multiple
representations would make
\[
 (N!+1)-N!=1
\]
divisible by $p$. This is impossible for an integer $p\geq2$:
if $1=pq$ for an integer $q$, then $q\ne0$ and
$1=p\abs q\geq p\geq2$ after taking absolute values, a
contradiction. Therefore $p>N>B$. This proves that the primes
are unbounded above.

Finally let $L>0$ be an arbitrary proposed lower bound for $P$.
Apply the fact just proved with $B=1/L$ to obtain a prime $p>1/L$.
Multiplication by $L>0$ gives $pL>1$, and division by $p>0$
then gives $1/p<L$. This member of $P$ defeats the proposed
lower bound $L$. Therefore
\[
 \inf P=0,\qquad \sup P=\frac12.
\]
Only the supremum is attained. Three upper bounds are $1/2,1,2$;
three lower bounds are $0,-1,-2$.

\phantomsection\label{bounds:t}
\step{Analysis (t): The real set $\{x\in\R:x^3<8\}$.}
We first identify the inequality's exact solution set. Multiplication
shows
\begin{align*}
 (x-2)(x^2+2x+4)
 &=x^3+2x^2+4x-2x^2-4x-8\\
 &=x^3-8.
\end{align*}
The second factor is always positive, because
\[
 x^2+2x+4=(x+1)^2+3\geq3>0.
\]
The sign of $x^3-8$ is therefore the sign of $x-2$. It follows
that $x^3<8$ if and only if $x<2$.

Thus $2$ is an upper bound. If $U<2$, the midpoint
$x=(U+2)/2$ satisfies
\[
 x-U=\frac{2-U}{2}>0,
 \qquad 2-x=\frac{2-U}{2}>0.
\]
Hence $U<x<2$, so $x$ belongs to the set and is above $U$.
This proves that the supremum is $2$. It is not attained since
$2^3=8$ fails the strict inequality. Upper bounds include $2,3,4$.

For any $L\in\R$, let $x=\min\{L,2\}-1$. Then $x<L$ and
$x<2$, so $x$ is a member below the proposed lower bound $L$.
Thus the set is NBB and has NO inf.

\phantomsection\label{bounds:u}
\step{Analysis (u): The set of real squares $\{x^2:x\in\R\}$.}
Every square is nonnegative, so $0$ is a lower bound. Taking
$x=0$ gives the member $0^2=0$. Therefore $0$ is the minimum
and infimum, and $0,-1,-2$ are lower bounds.

To defeat a proposed upper bound $M\in\R$, use Tool 1 to choose
$n\in\N$ with $n>\max\{M,1\}$. Since $n\geq1$,
$n^2\geq n>M$. Taking $x=n$ produces the member $x^2=n^2$
above $M$. Thus this set is NBA and has NO sup.

\phantomsection\label{bounds:v}
\step{Analysis (v): The values of $\cos(n\pi/3)$.}
We record a period table, which also supplies the proof for (w).
For a positive integer $n$, division of $n-1$ by $6$ gives
$n=6q+r$ with $q\geq0$ an integer and $r\in\{1,2,3,4,5,6\}$.
Therefore
\[
 \frac{n\pi}{3}=2q\pi+\frac{r\pi}{3}.
\]
The period $2\pi$ of sine and cosine shows that their values depend
only on $r$. The familiar unit-circle values are
\[
\begin{array}{c|c|c|c}
 r&r\pi/3&\cos(r\pi/3)&\sin(r\pi/3)\\ \hline
 1&\pi/3&1/2&\sqrt3/2\\
 2&2\pi/3&-1/2&\sqrt3/2\\
 3&\pi&-1&0\\
 4&4\pi/3&-1/2&-\sqrt3/2\\
 5&5\pi/3&1/2&-\sqrt3/2\\
 6&2\pi&1&0
\end{array}
\]
Every positive integer has one of these residues, so no other values
occur. Every row occurs at $n=r$, so none of the listed values is
missing from the corresponding set.

In particular the cosine set is exactly
\[
 \{-1,-1/2,1/2,1\}.
\]
Its minimum and infimum are $-1$, attained at $n=3$, and its
maximum and supremum are $1$, attained at $n=6$. Upper bounds
include $1,2,3$; lower bounds include $-1,-2,-3$.

\phantomsection\label{bounds:w}
\step{Analysis (w): The values of $\sin(n\pi/3)$.}
The period argument and all six sine values were displayed in
Analysis (v). They show that the sine set is exactly
\[
 \{-\sqrt3/2,0,\sqrt3/2\}.
\]
Because $\sqrt3>0$, these three numbers have the displayed
increasing order. Thus its infimum is the member $-\sqrt3/2$
at $n=4$, and its supremum is the member $\sqrt3/2$ at $n=1$.
Also $\sqrt3<2$: the contrary inequality would imply
$3=(\sqrt3)^2\geq4$. Hence $0<\sqrt3/2<1$, so all members
are between $-1$ and $1$. Three upper bounds are $1,2,3$;
three lower bounds are $-1,-2,-3$.

\subsubsection*{Exercise 4.1: Explicit answers for all twenty-three parts}

Each row gives three \emph{distinct} upper bounds, or NBA. A number
equal to the supremum is permitted as an upper bound whether or not
it belongs to the set. In a bounded row, the reason proves that the
smallest listed number is an upper bound; the larger two then follow
by transitivity of $\leq$. The linked analysis supplies the full
membership or unboundedness argument.

\begin{center}
\begin{tabular}{@{}c@{\quad}p{0.27\linewidth}@{\quad}p{0.59\linewidth}@{}}
\toprule
Part&Three upper bounds / NBA&Why these are valid\\ \midrule
(a)&$1,2,3$&Every member is at most $1$; \hyperref[bounds:a]{Analysis (a)}.\\
(b)&$1,2,3$&Every member is less than $1$; \hyperref[bounds:b]{Analysis (b)}.\\
(c)&$7,8,9$&Both members are at most $7$; \hyperref[bounds:c]{Analysis (c)}.\\
(d)&$4,5,6$&$e<3<\pi<4$; \hyperref[bounds:d]{Analysis (d)}.\\
(e)&$1,2,3$&$1/n\leq1$ for every $n\geq1$; \hyperref[bounds:e]{Analysis (e)}.\\
(f)&$0,1,2$&The only member is $0$; \hyperref[bounds:f]{Analysis (f)}.\\
(g)&$3,4,5$&Each interval lies at or below $3$; \hyperref[bounds:g]{Analysis (g)}.\\
(h)&NBA&The members $2n$ exceed any proposed bound; \hyperref[bounds:h]{Analysis (h)}.\\
(i)&$1,2,3$&The intersection equals $[0,1]$; \hyperref[bounds:i]{Analysis (i)}.\\
(j)&$1,2,3$&$1-3^{-n}<1$; \hyperref[bounds:j]{Analysis (j)}.\\
(k)&NBA&Members are at least $n-1$, unbounded above; \hyperref[bounds:k]{Analysis (k)}.\\
(l)&$2,3,4$&Membership requires $r<2$; \hyperref[bounds:l]{Analysis (l)}.\\
\bottomrule
\end{tabular}
\end{center}

\begin{center}
\begin{tabular}{@{}c@{\quad}p{0.27\linewidth}@{\quad}p{0.59\linewidth}@{}}
\toprule
Part&Three upper bounds / NBA&Why these are valid\\ \midrule
(m)&$2,3,4$&$r^2<4$ implies $r<2$; \hyperref[bounds:m]{Analysis (m)}.\\
(n)&$2,3,4$&$r<\sqrt2<2$; \hyperref[bounds:n]{Analysis (n)}.\\
(o)&$0,1,2$&Membership requires $x<0$; \hyperref[bounds:o]{Analysis (o)}.\\
(p)&$10,11,12$&$10$ is the largest of the four members; \hyperref[bounds:p]{Analysis (p)}.\\
(q)&$16,17,18$&$16$ is the largest listed member; \hyperref[bounds:q]{Analysis (q)}.\\
(r)&$1,2,3$&The intersection is $\{1\}$; \hyperref[bounds:r]{Analysis (r)}.\\
(s)&$1/2,1,2$&Every prime is at least $2$; \hyperref[bounds:s]{Analysis (s)}.\\
(t)&$2,3,4$&$x^3<8$ is equivalent to $x<2$; \hyperref[bounds:t]{Analysis (t)}.\\
(u)&NBA&Squares $n^2$ exceed any real bound; \hyperref[bounds:u]{Analysis (u)}.\\
(v)&$1,2,3$&The four cosine values are at most $1$; \hyperref[bounds:v]{Analysis (v)}.\\
(w)&$1,2,3$&The largest sine value is $\sqrt3/2<1$; \hyperref[bounds:w]{Analysis (w)}.\\
\bottomrule
\end{tabular}
\end{center}

\answer{The only sets in Exercise 4.1 that are not bounded above
are (h), (k), and (u). Every other part has the three distinct upper
bounds explicitly listed and justified above.}

\exercise{4.2}{Three lower bounds for each of the same sets}

\goal For each set (a)--(w) from Exercise 4.1, list three distinct
real lower bounds or give NBB, meaning ``not bounded below.''

\idea Use the lower-endpoint conclusions of the shared analyses.
If $\ell$ is a lower bound and $L\leq\ell$, then
$L\leq\ell\leq s$ for each member $s$, so $L$ is also a lower
bound. Consequently it suffices in each bounded row to justify the
largest of the three listed candidates; the other two are smaller.
For an NBB row, the linked proof constructs, for every $L\in\R$,
a member strictly below $L$.

\begin{center}
\begin{tabular}{@{}c@{\quad}p{0.27\linewidth}@{\quad}p{0.59\linewidth}@{}}
\toprule
Part&Three lower bounds / NBB&Why these are valid\\ \midrule
(a)&$0,-1,-2$&Every member is at least $0$; \hyperref[bounds:a]{Analysis (a)}.\\
(b)&$0,-1,-2$&Every member is positive; \hyperref[bounds:b]{Analysis (b)}.\\
(c)&$2,1,0$&Both members are at least $2$; \hyperref[bounds:c]{Analysis (c)}.\\
(d)&$2,1,0$&$2<e<\pi$; \hyperref[bounds:d]{Analysis (d)}.\\
(e)&$0,-1,-2$&Every $1/n$ is positive; \hyperref[bounds:e]{Analysis (e)}.\\
(f)&$0,-1,-2$&The only member is $0$; \hyperref[bounds:f]{Analysis (f)}.\\
(g)&$0,-1,-2$&Both intervals lie at or above $0$; \hyperref[bounds:g]{Analysis (g)}.\\
(h)&$2,1,0$&Each interval starts at $2n\geq2$; \hyperref[bounds:h]{Analysis (h)}.\\
(i)&$0,-1,-2$&The intersection is $[0,1]$; \hyperref[bounds:i]{Analysis (i)}.\\
(j)&$2/3,0,-1$&$1-3^{-n}\geq2/3$; \hyperref[bounds:j]{Analysis (j)}.\\
(k)&$0,-1,-2$&$n+(-1)^n/n\geq0$; \hyperref[bounds:k]{Analysis (k)}.\\
(l)&NBB&Negative integers in the set go below any $L$; \hyperref[bounds:l]{Analysis (l)}.\\
\bottomrule
\end{tabular}
\end{center}

\begin{center}
\begin{tabular}{@{}c@{\quad}p{0.27\linewidth}@{\quad}p{0.59\linewidth}@{}}
\toprule
Part&Three lower bounds / NBB&Why these are valid\\ \midrule
(m)&$-2,-3,-4$&$r^2<4$ implies $r>-2$; \hyperref[bounds:m]{Analysis (m)}.\\
(n)&$-2,-3,-4$&$r>-\sqrt2>-2$; \hyperref[bounds:n]{Analysis (n)}.\\
(o)&NBB&$\min\{L,0\}-1$ is a member below $L$; \hyperref[bounds:o]{Analysis (o)}.\\
(p)&$1,0,-1$&$1$ is the smallest of the four members; \hyperref[bounds:p]{Analysis (p)}.\\
(q)&$0,-1,-2$&$0$ is the smallest listed member; \hyperref[bounds:q]{Analysis (q)}.\\
(r)&$1,0,-1$&The only member of the intersection is $1$; \hyperref[bounds:r]{Analysis (r)}.\\
(s)&$0,-1,-2$&All prime reciprocals are positive; \hyperref[bounds:s]{Analysis (s)}.\\
(t)&NBB&$\min\{L,2\}-1$ satisfies $x^3<8$ and $x<L$; \hyperref[bounds:t]{Analysis (t)}.\\
(u)&$0,-1,-2$&Every real square is nonnegative; \hyperref[bounds:u]{Analysis (u)}.\\
(v)&$-1,-2,-3$&The cosine values are at least $-1$; \hyperref[bounds:v]{Analysis (v)}.\\
(w)&$-1,-2,-3$&The least sine value is $-\sqrt3/2>-1$; \hyperref[bounds:w]{Analysis (w)}.\\
\bottomrule
\end{tabular}
\end{center}

\answer{The only sets in Exercise 4.2 that are not bounded below
are (l), (o), and (t). Each remaining part has three distinct lower
bounds in its row.}

\exercise{4.3}{The real supremum of each set}

\goal For each of the twenty-three sets, give its supremum when
that real least upper bound exists; otherwise write NO sup.

\idea A mere upper bound is insufficient: we also need leastness.
The shared analyses prove leastness either by exhibiting the proposed
bound as a member or by constructing members above every smaller
candidate. The word ``attained'' in the table means the supremum
belongs to the set, in which case it is also the maximum. Failure
of attainment does not prevent a supremum from existing.

\begin{center}
\begin{tabular}{@{}c@{\quad}p{0.27\linewidth}@{\quad}p{0.59\linewidth}@{}}
\toprule
Part&Supremum&Bound and leastness proof\\ \midrule
(a)&$1$&An upper bound and a member; attained. \hyperref[bounds:a]{Analysis (a)}.\\
(b)&$1$&Upper endpoint is sharp by Tool 2; not attained. \hyperref[bounds:b]{Analysis (b)}.\\
(c)&$7$&The larger member is an upper bound; attained. \hyperref[bounds:c]{Analysis (c)}.\\
(d)&$\pi$&$e<\pi$ and $\pi$ is a member; attained. \hyperref[bounds:d]{Analysis (d)}.\\
(e)&$1$&$1/n\leq1$, with equality at $n=1$; attained. \hyperref[bounds:e]{Analysis (e)}.\\
(f)&$0$&The sole member is its upper bound; attained. \hyperref[bounds:f]{Analysis (f)}.\\
(g)&$3$&All members are at most $3$ and $3$ is a member; attained. \hyperref[bounds:g]{Analysis (g)}.\\
(h)&NO sup&There is no real upper bound at all. \hyperref[bounds:h]{Analysis (h)}.\\
(i)&$1$&The intersection equals $[0,1]$; attained. \hyperref[bounds:i]{Analysis (i)}.\\
(j)&$1$&Members exceed any $U<1$ but never equal $1$; not attained. \hyperref[bounds:j]{Analysis (j)}.\\
(k)&NO sup&The members exceed every real upper-bound candidate. \hyperref[bounds:k]{Analysis (k)}.\\
(l)&$2$&Density gives a rational member above any $U<2$; not attained. \hyperref[bounds:l]{Analysis (l)}.\\
\bottomrule
\end{tabular}
\end{center}

\begin{center}
\begin{tabular}{@{}c@{\quad}p{0.27\linewidth}@{\quad}p{0.59\linewidth}@{}}
\toprule
Part&Supremum&Bound and leastness proof\\ \midrule
(m)&$2$&The rational interval has sharp upper endpoint $2$; not attained. \hyperref[bounds:m]{Analysis (m)}.\\
(n)&$\sqrt2$&Density proves the upper endpoint is sharp; not attained. \hyperref[bounds:n]{Analysis (n)}.\\
(o)&$0$&For $U<0$, the member $U/2$ is above $U$; not attained. \hyperref[bounds:o]{Analysis (o)}.\\
(p)&$10$&All four values are at most the member $10$; attained. \hyperref[bounds:p]{Analysis (p)}.\\
(q)&$16$&The largest listed member is $16$; attained. \hyperref[bounds:q]{Analysis (q)}.\\
(r)&$1$&The intersection is the singleton $\{1\}$; attained. \hyperref[bounds:r]{Analysis (r)}.\\
(s)&$1/2$&Every prime is at least $2$, and $2$ is prime; attained. \hyperref[bounds:s]{Analysis (s)}.\\
(t)&$2$&The set is $\{x:x<2\}$, with members above every $U<2$; not attained. \hyperref[bounds:t]{Analysis (t)}.\\
(u)&NO sup&Real squares exceed every proposed upper bound. \hyperref[bounds:u]{Analysis (u)}.\\
(v)&$1$&The period table has maximum $1$, occurring at $n=6$; attained. \hyperref[bounds:v]{Analysis (v)}.\\
(w)&$\sqrt3/2$&The period table has this largest value at $n=1$; attained. \hyperref[bounds:w]{Analysis (w)}.\\
\bottomrule
\end{tabular}
\end{center}

\answer{Parts (h), (k), and (u) have NO sup in $\R$. Every other
supremum is the real number stated in its row, with its upper-bound
and leastness requirements proved in the linked analysis.}

\exercise{4.4}{The real infimum of each set}

\goal Give the real greatest lower bound of each set (a)--(w),
or write NO inf if the set has no real infimum.

\idea Reverse the roles of upper and lower bounds in Exercise 4.3.
For an attained lower bound, the member itself defeats every larger
candidate. For a nonattained infimum, the linked analysis constructs
members below every larger candidate. ``NO inf'' is used exactly
where there is no real lower bound.

\begin{center}
\begin{tabular}{@{}c@{\quad}p{0.27\linewidth}@{\quad}p{0.59\linewidth}@{}}
\toprule
Part&Infimum&Bound and greatestness proof\\ \midrule
(a)&$0$&A lower bound and a member; attained. \hyperref[bounds:a]{Analysis (a)}.\\
(b)&$0$&The open interval has sharp lower endpoint $0$; not attained. \hyperref[bounds:b]{Analysis (b)}.\\
(c)&$2$&The smaller member is a lower bound; attained. \hyperref[bounds:c]{Analysis (c)}.\\
(d)&$e$&$e<\pi$ and $e$ is a member; attained. \hyperref[bounds:d]{Analysis (d)}.\\
(e)&$0$&Positive reciprocals lie below any $L>0$; not attained. \hyperref[bounds:e]{Analysis (e)}.\\
(f)&$0$&The sole member is its lower bound; attained. \hyperref[bounds:f]{Analysis (f)}.\\
(g)&$0$&Every member is at least the member $0$; attained. \hyperref[bounds:g]{Analysis (g)}.\\
(h)&$2$&Every interval starts at least at $2$, and $2$ occurs; attained. \hyperref[bounds:h]{Analysis (h)}.\\
(i)&$0$&The intersection is $[0,1]$; attained. \hyperref[bounds:i]{Analysis (i)}.\\
(j)&$2/3$&Every value is at least $2/3$, with equality at $n=1$; attained. \hyperref[bounds:j]{Analysis (j)}.\\
(k)&$0$&Every value is nonnegative, with equality at $n=1$; attained. \hyperref[bounds:k]{Analysis (k)}.\\
(l)&NO inf&Negative integer members go below any proposed lower bound. \hyperref[bounds:l]{Analysis (l)}.\\
\bottomrule
\end{tabular}
\end{center}

\begin{center}
\begin{tabular}{@{}c@{\quad}p{0.27\linewidth}@{\quad}p{0.59\linewidth}@{}}
\toprule
Part&Infimum&Bound and greatestness proof\\ \midrule
(m)&$-2$&Density makes the strict lower endpoint sharp; not attained. \hyperref[bounds:m]{Analysis (m)}.\\
(n)&$-\sqrt2$&Density makes the strict lower endpoint sharp; not attained. \hyperref[bounds:n]{Analysis (n)}.\\
(o)&NO inf&Negative real members go below every real $L$. \hyperref[bounds:o]{Analysis (o)}.\\
(p)&$1$&The other three members are greater than the member $1$; attained. \hyperref[bounds:p]{Analysis (p)}.\\
(q)&$0$&The least of the listed members is $0$; attained. \hyperref[bounds:q]{Analysis (q)}.\\
(r)&$1$&The intersection equals $\{1\}$; attained. \hyperref[bounds:r]{Analysis (r)}.\\
(s)&$0$&Unbounded primes give a reciprocal below any $L>0$; not attained. \hyperref[bounds:s]{Analysis (s)}.\\
(t)&NO inf&The half-line $x<2$ has members below every real $L$. \hyperref[bounds:t]{Analysis (t)}.\\
(u)&$0$&Every square is nonnegative and $0^2=0$; attained. \hyperref[bounds:u]{Analysis (u)}.\\
(v)&$-1$&The smallest period-table value occurs at $n=3$; attained. \hyperref[bounds:v]{Analysis (v)}.\\
(w)&$-\sqrt3/2$&The smallest sine value occurs at $n=4$; attained. \hyperref[bounds:w]{Analysis (w)}.\\
\bottomrule
\end{tabular}
\end{center}

\answer{Parts (l), (o), and (t) have NO inf in $\R$. Every other
part has the real infimum in its row. In particular, (e) and (s)
have infimum $0$ without containing $0$, while the open-interval
intersection in (r) has both supremum and infimum equal to its
actual member $1$.}
